\paragraph{~TQFT 的定义。}\label{sub:tqft-definition}
对每一对\trExact\trLagrange\trSubmanifold{}$(L_0,L_1)$，选取一个\trRegular\trFloerDatum{}%
$(H_{L_0,L_1},J_{L_0,L_1})$。把\trFloerCochainGroup{}$CF^{\mathrm{pr}}(L_0,L_1)$ 定义为
一个 $\KK$-\trVectorSpace，使每个 $y\in\scr{C}(L_0,L_1)$ 恰有一个\trBasis{}向量。
这个群上的\trBoundaryOperator{}是
\begin{equation}\label{eq:boundary-operator}
\partial(y_1)=\sum_{y_0}\#\scr{M}^*_Z(y_0,y_1)\,y_0,
\end{equation}
其中 $\#$ 表示把\trIsolated{}点的个数按模 $2$ 计数。$HF^{\mathrm{pr}}(L_0,L_1)$ 就是
$(CF^{\mathrm{pr}}(L_0,L_1),\partial)$ 的\trCohomology。接下来，对每个带%
\trLagrange\trLabel{}$\{L_C\}$ 与\trStriplikeEnd{}的\trPointedBoundary\trRiemann%
\trSurface{}$S$，选取一个与前面那些\trFloerDatum{}\trCompatible{}的%
\trRegular\trPerturbationDatum{}$(K_S,J_S)$。定义映射
\begin{equation}\label{eq:cphi-map}
\begin{aligned}
&C\Phi_S\colon
\bigotimes_{\zeta^+\in\Sigma^+}CF^{\mathrm{pr}}(L_{\zeta+,0},L_{\zeta+,1})
\longrightarrow
\bigotimes_{\zeta^-\in\Sigma^-}CF^{\mathrm{pr}}(L_{\zeta-,0},L_{\zeta-,1}),\\
&C\Phi_S(\otimes_{\zeta+}y_{\zeta+})
=\sum_{\{y_{\zeta-}\}}\#\scr{M}_S(\{y_{\zeta^-},y_{\zeta^+}\})\,(\otimes_{\zeta-}y_{\zeta^-}).
\end{aligned}
\end{equation}
并取 $\Phi_S$ 为在\trCohomology{}上诱导的映射。非专业的读者应当记住：尽管这里的处理
相当粗略，这项构造的每一块都是一件浩大的工程。光是定义 $\partial$ 与 $C\Phi_S$，
就需要 (i) \trTransversality{}结果——它们说，对\trFloerDatum{}与%
\trPerturbationDatum{}而言，\trRegularity{}是一个\trGeneric{}条件；以及
(ii) \trCompactness{}结果——它们保证在\trRegular{}的情形，相关的\trModuliSpace{}%
只有有限多个\trIsolated{}点。在那之后，一些适当的\trGluing{}定理表明 (iii)
$\partial^2=0$，以及 (iv) 诸 $C\Phi_S$ 都是\trChainMap。进一步的步骤是：(v) 所得的
理论与一切选取无关（至多相差\trCanonical\trIsomorphism）；以及 (vi) 它满足%
\trComposition{}公理\eqref{eq:composition-law}。对前四个问题我们将一字不提，因为
文献对它们的论述已相当完备，而且在构造\trFukayaCategory{}（小节\cref{sub:transversality} 与\cref{sub:compactness}）时
它们反正还会出现。关于最后两步（它们实际上是相互交织的），几句评论或许是合适的。
为展开 (v)，人们首先证明诸 $\Phi_S$ 不依赖于 $S$ 的\trConformal{}结构、诸%
\trStriplikeEnd{}的选取以及\trPerturbationDatum，只要诸端处的\trFloerDatum{}保持
固定。为此，人们取一族 $[0;1]$ 上的\trConformal{}结构与\trPerturbationDatum，在两组
给定的之间\trInterpolation，并利用由此所得的全部\trModuliSpace{}之并来定义一个%
\trChainHomotopy。（这可以看作\trParametrization\trModuliSpace{}的第一次、而且相当
平凡的现身，我们稍后将对它们做详尽的讨论。）接下来转向 (vi)：取配备了%
\trStriplikeEnd、\trLagrange\trBoundaryCondition{}与\trPerturbationDatum{}的%
\trSurface{}$S_1$、$S_2$。我们想沿无穷远点 $\zeta_1^+$、$\zeta_2^-$ 把它们%
\trGluing{}起来以构成 $S=S_1\#_l S_2$，如\eqref{eq:striplike-gluing} 中那样。
假设与这两个端 $\zeta_1^+$、$\zeta_2^-$ 相关联的\trFloerDatum{}相重合，这样便自然地
得到 $S$ 上的一个\trPerturbationDatum。再假设适当的\trRegularity{}条件成立，并且
$l\gg0$，则可以证明\trCompositionLaw\emph{在链的层面}上成立。当然，这个等式只对
$S$ 上\trConformal{}结构与\trPerturbationDatum{}的这一非常特定的选取成立。然而，我们
已经证明了\trCohomology{}层面的映射 $\Phi_S$ 与这些选取无关。（对于%
\trFukayaCategory——它们是链层面的结构——这类论证将不再适用，这意味着我们将被迫对
构造中所用的\trPerturbationDatum{}施加多得多的控制。）鉴于此，(v) 的剩余部分——
它关乎\trFloerDatum{}选取的无关性——就完全是形式上的了。考虑无限\trStrip{}$Z$，
在其边界\trConnectedComponent{}$C_k=\RR\times\{k\}$（$k=0,1$）上贴有%
\trLagrange\trLabel{}$L_k$。在两个端 $\zeta_0=-\infty$、$\zeta_1=+\infty$ 处，给对
$(L_{\zeta,0},L_{\zeta,1})=(L_0,L_1)$ 配备\trRegular\trFloerDatum{}%
$(H_{\zeta_j},J_{\zeta_j})=(H^j_{L_0,L_1},J^j_{L_0,L_1})$，并选取一个%
\trRegular{}而\trCompatible{}的\trPerturbationDatum。这就给出了一个映射
\begin{equation*}
\Phi_Z\colon HF^{\mathrm{pr}}(L_0,L_1)^1\longrightarrow HF^{\mathrm{pr}}(L_0,L_1)^0
\end{equation*}
介于由此所得的诸\trFloerCohomology{}群之间，传统上称为\emph{\trContinuationMap}
（continuation map）[118]。上文提到的\trHomotopy{}论证表明它是\trCanonical{}的；
而\trCompositionLaw{}的一个应用表明它是一个\trIsomorphism。
\begin{remark}\label{thm:vanishing-inhomogeneous-terms}
在许多情形，可以使用\trInhomogeneous{}项消没的\trFloerDatum{}与%
\trPerturbationDatum。也就是说，每当 $L_0\pitchfork L_1$，就能找到一个 $H=0$ 的%
\trRegular\trFloerDatum。类似地，考虑一个带\trStriplikeEnd{}与%
\trLagrange\trBoundaryCondition{}$\{L_C\}$ 的\trRiemann\trSurface{}$S$。假设与无穷远点
相关联的任何一对\trLagrange\trSubmanifold{}都\trTransverse{}地相交，并且 $M$ 中
没有同时落在全部 $L_C$ 上的点；第二个条件排除了\eqref{eq:inhomogeneous-equation}
的常值解。于是就能找到一个 $K=0$ 的\trRegular\trPerturbationDatum。
\end{remark}
